\section{Modelado del comportamiento y fundamentos del aprendizaje por refuerzo}

\subsection{Representación de datos de la experiencia}

\begin{importantbox}{Términos clave}
\begin{itemize}
    \item \textbf{Estado} $s_t$: el estado del mundo en el instante $t$.
    \item \textbf{Observación} $o_t$: la información que el agente observa en el instante $t$ (se usa cuando hay pérdida de información).
    \item \textbf{Acción} $a_t$: la decisión en el instante $t$.
    \item \textbf{Trayectoria} $\tau = (s_1, a_1, s_2, a_2, \ldots, s_T, a_T)$: una secuencia completa de estados y acciones.
    \item \textbf{Función de recompensa} $r(s, a)$: mide qué tan buena es una pareja estado-acción.
\end{itemize}
\end{importantbox}

\subsection{Diferencia entre estado y observación}

Una propiedad clave es la \textbf{propiedad de Markov} (Markov Property): el siguiente estado depende únicamente del estado y la acción actuales, y no de la historia anterior:
$$p(s_{t+1} \mid s_t, a_t) \quad \text{independiente de } s_{t-1}$$

Pero en muchos escenarios reales (como los chatbots), el agente solo puede ver la \textbf{observación} y no el estado completo, por lo que hay que agregar memoria a la política: $\pi_\theta(a_t \mid o_{t-m}, \ldots, o_t)$.

\subsection{Representar el comportamiento con una red neuronal}

La política $\pi_\theta(a \mid s)$ es la función que mapea estados a una distribución sobre acciones. El proceso de ejecución del agente es el siguiente:
\begin{enumerate}
    \item Observa el estado $s_t$
    \item Muestrea una acción de la política $a_t \sim \pi_\theta(\cdot \mid s_t)$
    \item Observa el siguiente estado $s_{t+1} \sim p(\cdot \mid s_t, a_t)$
\end{enumerate}
La secuencia completa que se produce de esta manera se llama rollout de la política o episodio.

\subsection{El objetivo del aprendizaje por refuerzo}

El objetivo del aprendizaje por refuerzo es \textbf{maximizar la recompensa acumulada esperada}:
$$\max_\theta \; \mathbb{E}_{\tau \sim p_\theta(\tau)} \left[ \sum_{t} r(s_t, a_t) \right]$$

donde la distribución sobre trayectorias es:
$$p_\theta(\tau) = p(s_1) \prod_{t=1}^{T} \pi_\theta(a_t \mid s_t) \, p(s_{t+1} \mid s_t, a_t)$$

\begin{warningbox}{Por qué usar una esperanza}
La recompensa acumulada no es una cantidad determinista; tiene dos fuentes de aleatoriedad: (1) el entorno mismo es aleatorio; (2) la política puede ser aleatoria. Por eso lo que optimizamos es la recompensa acumulada \textbf{esperada}.
\end{warningbox}

\subsection{El sentido de las políticas aleatorias}

¿Por qué usar una política aleatoria en lugar de una política determinista? Por dos razones:
\begin{enumerate}
    \item \textbf{Exploración}: para aprender de la propia experiencia, hay que probar cosas distintas.
    \item \textbf{Modelar comportamiento aleatorio}: los datos existentes suelen mostrar comportamientos diversos, y se pueden aprovechar las herramientas de los modelos generativos.
\end{enumerate}

\subsection{La función de valor y la función Q}

Dos cantidades centrales para medir qué tan buena es una política:

\begin{itemize}
    \item \textbf{Función de valor} $V^\pi(s)$: la recompensa futura esperada al partir del estado $s$ y seguir la política $\pi$.
    \item \textbf{Función Q} $Q^\pi(s, a)$: la recompensa futura esperada al tomar la acción $a$ desde el estado $s$ y luego seguir la política $\pi$.
\end{itemize}

\subsection{Panorama de los tipos de algoritmos}

\begin{enumerate}
    \item \textbf{Imitation Learning}: imita una política capaz de obtener una recompensa alta.
    \item \textbf{Policy Gradients}: calcula directamente el gradiente de la función objetivo.
    \item \textbf{Actor-Critic}: estima el valor de la política actual y lo usa para mejorar la política.
    \item \textbf{Value-based}: estima el valor de la política óptima.
    \item \textbf{Model-based}: aprende un modelo de la dinámica del entorno, para usarlo en la planeación o en la mejora de la política.
\end{enumerate}

Los distintos algoritmos hacen distintas concesiones, y son adecuados para distintas condiciones supuestas (costo de recolección de datos, forma de la supervisión, estabilidad, características del espacio de acciones, etc.).

\subsection{Resumen de la sección}

Los elementos básicos del aprendizaje por refuerzo incluyen el estado, la acción, la trayectoria, la recompensa y la política. El objetivo es maximizar la recompensa acumulada esperada. La función de valor y la función Q son las herramientas centrales para evaluar qué tan buena es una política. Los distintos tipos de algoritmos tienen énfasis diferentes y son adecuados para escenarios diferentes.
